\documentclass[exercice]{thedocument}%
\usepackage{tkz-euclide}
\startdatakeys
    \source{18GENMATMEAG1}
    \BO{2026}
    \cycle{4}
    \competencesDuSocle{RA,CO}
    \objectifsApprentissage{B3ca}
\enddatakeys
\begin{document}%
    La figure ci-dessous n’est pas réalisée en vraie grandeur.
    \begin{center}
        \begin{tikzpicture}[scale=0.8]
            \coordinate[label=above:\mathspoint{A}](A)at(0,0);
            \coordinate[label=above right:\mathspoint{B}](B)at(160:4.5);
            \coordinate[label=below left:\mathspoint{C}](C)at(195:3);
            \coordinate[label=below:\mathspoint{E}](E)at(15:5);
            \coordinate[label=below left:\mathspoint{F}](F)at(-20:7.5);
            \coordinate[label=below left:\mathspoint{N}](N)at(195:4.8);
            \coordinate[label=below left:\mathspoint{M}](M)at(160:7.2);
            \foreach\X in{A,B,C,E,F,N,M}{
                \draw node at(\X){$\times$};
                }
            \tkzDrawLines[add=0.2 and 0.2](M,F N,E B,C M,N)
        \end{tikzpicture}
    \end{center}
    Les droites \mathspoint{(BC)} et \mathspoint{(MN)} sont
    parallèles.\medskip\par
    \begin{enumerate}
        \item Calculer \mathspoint{AM} et \mathspoint{BC}.
        \item On sait de plus que $\mathspoint{AE}=5$~cm  et $\mathspoint{AF}=
        7,5$~cm.\smallskip\par
        Montrer que les droites \mathspoint{(EF)} et \mathspoint{(BC)} sont
        parallèles.
    \end{enumerate}
\end{document}